%%%%%%%%%%%%%%%%%%%%%%%%%%%%%%%%%%%%%%%%%%%%%%%%%%%%%%%%%%%%%%%%%%%%%%%%%%%%%%%%
%        Reviewer Response — UL Report
%        CS7641 Machine Learning — Spring 2026
%%%%%%%%%%%%%%%%%%%%%%%%%%%%%%%%%%%%%%%%%%%%%%%%%%%%%%%%%%%%%%%%%%%%%%%%%%%%%%%%

\documentclass[letterpaper, 10.5pt]{article}

\usepackage[margin=0.75in, top=0.7in, bottom=0.7in]{geometry}
\usepackage[utf8]{inputenc}
\usepackage[T1]{fontenc}
\usepackage{booktabs}
\usepackage{hyperref}
\usepackage{xcolor}
\usepackage{soul}
\sethlcolor{yellow!40}
\usepackage{parskip}
\setlength{\parskip}{3pt}
\setlength{\parindent}{0pt}
\usepackage{titlesec}
\titlespacing*{\section}{0pt}{6pt}{3pt}
\usepackage{enumitem}
\setlist[itemize]{topsep=1pt, itemsep=0pt, parsep=0pt, leftmargin=1.5em}

\title{\vspace{-1.5em}Reviewer Response --- UL Report\\[3pt]
\large Unsupervised Learning and Dimensionality Reduction\\
on Adult Income and Wine Quality\vspace{-0.5em}}

\author{Siddartha Sagar Chinne \quad \texttt{schinne3@gatech.edu}}
\date{April 26, 2026\vspace{-1em}}

\begin{document}

\maketitle
\thispagestyle{empty}
\pagestyle{empty}

\noindent\textbf{How to find the changes.}
All new or corrected text, new figures, and updated captions in the revised
report are \hl{highlighted in light yellow}.  New figure captions (Figs.~7
and~8) are highlighted in full.

%%%%%%%%%%%%%%%%%%%%%%%%%%%%%%%%%%%%%%%%%%%%%%%%%%%%%%%%%%%%%%%%%%%%%%%%%%%%%%%%
\section*{A3 --- Graph/Text Legibility (--2 pts)}

\textbf{Comment.} Legends, axis labels, and plot values were too small to
verify analysis without zoom.

\textbf{Root cause.} No global font-size override was applied to matplotlib;
figures exported at 8--9\,pt scaled below the legibility threshold in the
two-column IEEE layout.

\textbf{Fix.} A global font-size override was applied to the plotting
pipeline (axis labels 11\,pt, titles 12\,pt, tick labels 10\,pt) and all
figures were regenerated.  Figs.~7--8 (new K-selection sweep figures)
were also exported at a figsize calibrated to render at near 1:1 scale
within the column width, so each of the six panels displays at the
correct landscape aspect ratio with readable axis labels.  Every figure
in the revised report is readable at normal zoom without relying on
magnification.

%%%%%%%%%%%%%%%%%%%%%%%%%%%%%%%%%%%%%%%%%%%%%%%%%%%%%%%%%%%%%%%%%%%%%%%%%%%%%%%%
\section*{E --- Step 3: DR then Clustering (--8 pts)}

\textbf{Comment.} Unclear how cluster count was selected in reduced spaces.
Silhouette and AIC/BIC plots needed to justify selection.
Orderly visualization of all six DR--Clustering combinations missing for
both datasets.

\textbf{Root cause.} Step~3 reused raw-space frozen K values ($K\!=\!2/8$,
$N\!=\!7$) in all six reduced spaces.  The internal spec said ``no
re-selection,'' diverging from the assignment intent.  No reduced-space
sweep was ever run, so the report could show final results but not the
selection evidence.

\textbf{Fix --- K re-selection.}
The experiment was redesigned so that each of the six reduced spaces
(Wine and Adult $\times$ PCA/ICA/RP) gets its own independent clustering
sweep over $K\in[2,20]$.  KMeans selects K by highest silhouette (CH
tiebreaker, DB secondary) and GMM selects by lowest BIC --- identical
criteria to Step~1.  The resulting reduced-space K values are materially
different from the raw-space frozen values and are now justified by
visible sweep evidence.

\textbf{Fix --- new figures (captions fully highlighted).}
\begin{itemize}
  \item \textbf{Fig.~7 (new):} Wine reduced-space K selection sweeps ---
        3 rows (PCA/ICA/RP) $\times$ 2 cols (KMeans silhouette, GMM
        BIC+AIC).  Dashed vertical line marks selected K per panel.
  \item \textbf{Fig.~8 (new):} Adult reduced-space K selection sweeps,
        same layout.
\end{itemize}

The reduced-space K values differ materially from the raw-space frozen
values: Wine PCA KMeans retains $K\!=\!2$ (red/white split preserved), Wine
ICA KMeans selects $K\!=\!4$ (ICA exposes finer chemical structure), and
Adult GMM selects $\geq\!19$ components across all DR methods.

\textbf{Fix --- prose (highlighted in report).}
A \textbf{K re-selection} paragraph was added at the start of the Step~3
section.  ICA Wine sentences were updated to reference re-selected
$K\!=\!4$.  A factual error was also corrected: the original stated ``both
methods substantially improve Adult GMM'' --- PCA in fact \emph{degrades}
Adult GMM silhouette by $-17\%$ and that sentence now says so explicitly.

%%%%%%%%%%%%%%%%%%%%%%%%%%%%%%%%%%%%%%%%%%%%%%%%%%%%%%%%%%%%%%%%%%%%%%%%%%%%%%%%
\section*{G --- Step 5: Clustering on NN (--6 pts)}

\textbf{Comment.} Wall-clock analysis missing for EM and KMeans
cluster-feature variants.  Should relate runtime to the baseline NN and
mention the filtering effect.

\textbf{Root cause.} \texttt{run\_phase\_6\_nn\_cluster\_features.py}
stored only F1 metrics.  No timing instrumentation existed.

\textbf{Fix --- timing data.}
Wall-clock time was measured per seed for all three cluster-feature
variants and the mean is now reported in the Step~5 results table
(new \textbf{Time (s)} column visible in the revised report).
Measured times: KMeans one-hot 0.28\,s, KMeans distances 0.24\,s,
GMM posterior 0.25\,s (mean over 10 seeds).

\textbf{Fix --- new Speed paragraph (fully highlighted in report).}
\begin{itemize}
  \item Cluster features are \emph{appended} to the raw 12-d input, so
        input dimension grows; no speedup is expected or observed.
  \item Cluster-feature variants (0.24--0.28\,s) match the Step~4
        DR-reduced variants (0.24--0.28\,s) --- appending a few extra
        dimensions to a 12-d input carries negligible per-epoch overhead.
        The performance gain is from the filtering/routing signal
        the network gains for free, not from reduced arithmetic cost.
  \item The Step~4 raw baseline (0.70\,s) is not directly comparable
        because it includes PyTorch's one-time JIT compilation overhead;
        after warmup, per-epoch cost differences are $\sim$2--3\,ms.
\end{itemize}

\end{document}
